\section{Method}
\label{sec:method}

We consider knowledge distillation for instruction-conditioned multimodal
retrieval. A frozen teacher and a trainable student may use different
architectures and embedding dimensions. Our objective therefore transfers the
teacher's relational structure without aligning embedding coordinates. The
central object is the correspondence-aware, multiscale connectivity hierarchy
induced by query--candidate distances: it records both when vertices become
connected and which query and candidate identities participate in each
connection. Throughout this section, identity refers to row-wise
teacher--student correspondence, not to an additional dataset class-label
annotation.

\subsection{Problem Setup}
\label{sec:method-setup}

Let a global training batch contain $B$ query--positive pairs
$\mathcal{B}=\{(x_i,y_i)\}_{i=1}^{B}$ drawn from a single retrieval task. For
model $M\in\{\teacher,\student\}$, we denote its $\ell_2$-normalised query and
positive embeddings by
\begin{equation}
    \bm{q}^{M}_i
    = \frac{f_M(x_i)}{\lVert f_M(x_i)\rVert_2},
    \qquad
    \bm{p}^{M}_i
    = \frac{f_M(y_i)}{\lVert f_M(y_i)\rVert_2}.
    \label{eq:normalised-embeddings}
\end{equation}
The teacher and student dimensions, $d_{\teacher}$ and $d_{\student}$, need
not be equal. The student is trained with the standard in-batch retrieval loss
\begin{equation}
    \lossret
    = -\frac{1}{B}\sum_{i=1}^{B}
    \log
    \frac{
        \exp\!\left((\bm{q}^{\student}_i)^{\top}
        \bm{p}^{\student}_i / \tau\right)
    }{
        \sum_{j=1}^{B}
        \exp\!\left((\bm{q}^{\student}_i)^{\top}
        \bm{p}^{\student}_j / \tau\right)
    },
    \label{eq:retrieval-loss}
\end{equation}
where $\tau$ is the temperature.

Some datasets contain several queries with the same target, as is common when
class labels or short answers form the candidate set. Treating repeated
targets as distinct graph vertices would make the topology depend on batch
multiplicity rather than on the retrieval problem. We therefore assign every
target a deterministic, task-local identity and retain one representative for
each identity. This produces $C\leq B$ canonical candidates
$\{c_j\}_{j=1}^{C}$, with normalised embeddings $\bm{c}^M_j$. All $B$ query
vertices are retained. Candidate
canonicalisation is used only to construct the structural target; the
retrieval loss in Equation~\eqref{eq:retrieval-loss} remains defined on the
original training rows.

\subsection{Cross-Modal Filtration}
\label{sec:method-filtration}

For each model, we compute the rectangular cosine-distance matrix between the
queries and canonical candidates,
\begin{equation}
    \distmat^{M}_{ij}
    = 1-(\bm{q}^{M}_i)^{\top}\bm{c}^{M}_j,
    \qquad
    \distmat^{M}\in[0,2]^{B\times C}.
    \label{eq:cross-modal-distance}
\end{equation}
Because the distances are computed separately in each representation space,
Equation~\eqref{eq:cross-modal-distance} does not require a
teacher-to-student projector.

We interpret $\distmat^{M}$ as a complete weighted bipartite graph with the
shared identity-aligned vertex set
$\vertices=\queries\sqcup\candidates$, where
$\queries=\{q_i\}_{i=1}^{B}$ and
$\candidates=\{c_j\}_{j=1}^{C}$. At threshold $\varepsilon$, its filtration is
\begin{equation}
    \mathcal{G}^{M}_{\varepsilon}
    = \left(
        \vertices,
        \edges^{M}_{\varepsilon}
      \right),
    \qquad
    \edges^{M}_{\varepsilon}
    = \left\{(q_i,c_j):
      \distmat^{M}_{ij}\leq\varepsilon\right\}.
    \label{eq:bipartite-filtration}
\end{equation}
The graph contains no query--query or candidate--candidate edges. Consequently,
any connectivity within either side is induced exclusively by the
cross-modal relation, for example through candidates shared by two queries.

\subsection{Correspondence-Aware Persistent Connectivity}
\label{sec:method-merge-times}

For any two indexed vertices $a,b\in\vertices$, define their merge time as
the first filtration threshold at which they belong to the same connected
component:
\begin{align}
    \mergemat^{M}_{ab}
    &= \inf\left\{\varepsilon:
       a\text{ and }b\text{ are connected in }
       \mathcal{G}^{M}_{\varepsilon}\right\} \\
    &= \min_{\pi:a\rightsquigarrow b}
       \max_{e\in\pi}\distmat^{M}_{e},
    \label{eq:merge-time}
\end{align}
where $\pi$ ranges over all alternating paths between $a$ and $b$. Thus,
$\mergemat^{M}$ is the cophenetic, or minimax, ultrametric of the cross-modal
graph. Its three blocks have direct interpretations:
$\mergemat_{q_i c_j}$ measures when a query joins a candidate, while
$\mergemat_{q_i q_k}$ and $\mergemat_{c_j c_\ell}$ measure connectivity through
shared vertices on the opposite side.

The complete matrix can be recovered exactly from any minimum spanning tree
(MST) of the weighted bipartite graph. In particular, for every pair $(a,b)$,
$\mergemat^{M}_{ab}$ equals the largest edge weight on the unique MST path from
$a$ to $b$. We exploit this property rather than explicitly sweeping over
filtration thresholds.

The distinction between this identity-preserving matrix and an ordinary $H_0$
persistence barcode is essential. The finite $H_0$ barcode contains only the
sorted multiset of the $B+C-1$ MST edge weights. It discards the identities of
the vertices merged by those edges. For example, let $P$ be a non-identity
permutation of candidate columns and suppose
$\distmat^{\student}=\distmat^{\teacher}P$. The resulting graphs are
isomorphic after permuting candidates and hence have identical $H_0$
barcodes. Nevertheless, under the fixed teacher--student correspondence, the
correct candidate attached to each query has changed. The entries of
$\mergemat^{\student}$ detect this change.

\subsection{Structural Distillation Objective}
\label{sec:method-objective}

Teacher and student share the vertex identities and ordering of $\vertices$,
so their merge-time matrices can be compared entry by entry. Let
$N=B+C$ and $\Omega=\{(a,b):1\leq a<b\leq N\}$. We define
\begin{equation}
    \losstopo
    = \frac{1}{|\Omega|}
      \sum_{(a,b)\in\Omega}
      \left|
        \mergemat^{\teacher}_{ab}
        - \mergemat^{\student}_{ab}
      \right|.
    \label{eq:topology-loss}
\end{equation}
The default objective uses all three blocks of the merge-time matrix. A
query--candidate-only variant, obtained by setting
$\Omega=\queries\times\candidates$, serves as an ablation of the connectivity
induced within each side.

Equation~\eqref{eq:topology-loss} also admits a direct multiscale
interpretation. Since every cosine distance lies in $[0,2]$, for any vertex
pair,
\begin{equation}
    \left|\mergemat^{\teacher}_{ab}-\mergemat^{\student}_{ab}\right|
    = \int_{0}^{2}
      \left|
      \mathbb{1}\!\left[\mergemat^{\teacher}_{ab}\leq\varepsilon\right]
      -
      \mathbb{1}\!\left[\mergemat^{\student}_{ab}\leq\varepsilon\right]
      \right|\,\mathrm{d}\varepsilon.
    \label{eq:multiscale-interpretation}
\end{equation}
Hence the loss is the average integrated disagreement between the
identity-aligned teacher and student connectivity partitions across all
filtration scales.

The final training objective contains one structural coefficient:
\begin{equation}
    \mathcal{L}
    = \lossret + \lambda_{\mathrm{topo}}\losstopo.
    \label{eq:full-objective}
\end{equation}
Only the teacher's final normalised embeddings are required. The teacher
remains frozen, and the method uses neither hidden states nor attention maps
nor a learned projector.

\subsection{Optimisation and Distributed Batching}
\label{sec:method-optimisation}

For each model, we first extract an MST from the detached distance matrix. A
Kruskal pass over the MST assigns a bottleneck witness edge to every indexed
vertex pair: whenever an edge joins components $A$ and $B$, that edge witnesses
the merge time of all pairs in $A\times B$. Each entry of $\mergemat^M$ is then
gathered from its witness in the original distance matrix. We treat the
discrete MST routing as fixed during a backward pass, while gradients flow
through the selected student distances. This gives the exact gradient almost
everywhere. Distance ties can change the MST combinatorics, and a zero residual
in the absolute-error objective introduces the usual subgradient ambiguity.

In distributed training, embeddings and their identities are gathered across
workers before constructing the graph, so the topology is defined on the
global rather than per-device batch. Each global batch is restricted to one
task, all query rows are preserved, and repeated candidate identities are
collapsed as described in Section~\ref{sec:method-setup}. These constraints
prevent unrelated task spaces and duplicate target vertices from changing the
structural supervision.
